% --------------------------------------------------
% Chapter 07
% Roles and Responsibilities
% --------------------------------------------------

\label{chap:roles-responsibilities}

\section{Learning Objectives}

After completing this chapter, learners will be able to:

\begin{itemize}
    \item Identify the key roles involved in MLOps.
    \item Understand responsibilities throughout the ML lifecycle.
    \item Explain interactions between teams.
    \item Build a RACI matrix for ML projects.
    \item Design an effective operating model.
\end{itemize}

\section{Introduction}

Successful machine learning systems require more than algorithms and infrastructure.

Production-grade AI systems depend on collaboration between multiple teams possessing complementary skills.

Historically, organizations often assumed that Data Scientists could independently deliver machine learning solutions from development to production.

Experience has demonstrated that this approach rarely scales.

Modern MLOps environments rely on specialized roles working together throughout the model lifecycle.

\section{Overview of Roles}

A typical MLOps organization includes:

\begin{itemize}
    \item Business Owner
    \item Product Owner
    \item Data Engineer
    \item Data Scientist
    \item Machine Learning Engineer
    \item DevOps Engineer
    \item Model Validator
    \item Risk Manager
    \item IT Security Officer
    \item Operations Team
\end{itemize}

Although responsibilities vary across organizations, these roles cover most enterprise machine learning environments.

\section{Business Owner}

\subsection{Purpose}

The Business Owner is accountable for the business value generated by the machine learning solution.

The Business Owner defines:

\begin{itemize}
    \item Business objectives
    \item Success criteria
    \item Budget priorities
    \item Strategic direction
\end{itemize}

\subsection{Typical Responsibilities}

Examples include:

\begin{itemize}
    \item Defining business requirements
    \item Approving project scope
    \item Evaluating business impact
    \item Supporting adoption
\end{itemize}

\subsection{Credit Scoring Example}

For a credit scoring platform, the Business Owner may be:

\begin{itemize}
    \item Head of Retail Banking
    \item Chief Risk Officer
    \item Head of Credit Risk
\end{itemize}

\section{Product Owner}

\subsection{Purpose}

The Product Owner acts as the bridge between business stakeholders and technical teams.

The Product Owner ensures that development activities remain aligned with business priorities.

\subsection{Responsibilities}

Typical responsibilities include:

\begin{itemize}
    \item Requirements management
    \item Backlog prioritization
    \item Stakeholder communication
    \item Sprint planning support
\end{itemize}

\subsection{Key Skills}

A Product Owner requires:

\begin{itemize}
    \item Business knowledge
    \item Communication skills
    \item Project management capabilities
    \item Basic technical understanding
\end{itemize}

\section{Data Engineer}

\subsection{Purpose}

Data Engineers build and maintain the data infrastructure required by machine learning systems.

Their work provides reliable access to high-quality data.

\subsection{Responsibilities}

Typical activities include:

\begin{itemize}
    \item Data ingestion
    \item Data transformation
    \item ETL development
    \item Data quality controls
    \item Data pipeline maintenance
\end{itemize}

\subsection{Key Technologies}

Examples include:

\begin{itemize}
    \item SQL
    \item Spark
    \item Airflow
    \item Databricks
    \item Azure Data Factory
\end{itemize}

\subsection{Success Criteria}

Data Engineers focus on:

\begin{itemize}
    \item Reliability
    \item Scalability
    \item Data quality
    \item Pipeline performance
\end{itemize}

\section{Data Scientist}

\subsection{Purpose}

Data Scientists are responsible for developing machine learning models that address business problems.

They transform data into predictive insights and analytical solutions.

\subsection{Responsibilities}

Typical activities include:

\begin{itemize}
    \item Exploratory Data Analysis (EDA)
    \item Feature engineering
    \item Model development
    \item Hyperparameter tuning
    \item Performance evaluation
    \item Model documentation
\end{itemize}

\subsection{Key Technologies}

Examples include:

\begin{itemize}
    \item Python
    \item Pandas
    \item NumPy
    \item Scikit-learn
    \item XGBoost
    \item Jupyter Notebooks
\end{itemize}

\subsection{Success Criteria}

Data Scientists focus on:

\begin{itemize}
    \item Model performance
    \item Explainability
    \item Robustness
    \item Business relevance
\end{itemize}

\subsection{Common Misconception}

Data Scientists should not be solely responsible for deploying and operating production systems.

Operationalization requires additional engineering expertise.

\section{Machine Learning Engineer}

\subsection{Purpose}

Machine Learning Engineers bridge the gap between data science and production operations.

They transform experimental models into reliable production services.

\subsection{Responsibilities}

Typical activities include:

\begin{itemize}
    \item Building training pipelines
    \item Model packaging
    \item Model deployment
    \item Experiment tracking integration
    \item Model registry management
    \item Monitoring implementation
\end{itemize}

\subsection{Key Technologies}

Examples include:

\begin{itemize}
    \item MLflow
    \item Docker
    \item Kubernetes
    \item Azure Machine Learning
    \item Kubeflow
    \item GitHub Actions
\end{itemize}

\subsection{Success Criteria}

ML Engineers focus on:

\begin{itemize}
    \item Reproducibility
    \item Automation
    \item Scalability
    \item Reliability
\end{itemize}

\subsection{Position Within the Organization}

ML Engineers often collaborate closely with:

\begin{itemize}
    \item Data Scientists
    \item Data Engineers
    \item DevOps Engineers
\end{itemize}

They are frequently considered the central operational role within MLOps initiatives.

\section{DevOps Engineer}

\subsection{Purpose}

DevOps Engineers provide the infrastructure and automation capabilities required to operate machine learning systems at scale.

Their expertise originates from software engineering and cloud operations.

\subsection{Responsibilities}

Typical responsibilities include:

\begin{itemize}
    \item Infrastructure provisioning
    \item CI/CD pipeline implementation
    \item Container orchestration
    \item Platform monitoring
    \item Security controls
    \item Capacity planning
\end{itemize}

\subsection{Key Technologies}

Examples include:

\begin{itemize}
    \item GitHub Actions
    \item Azure DevOps
    \item Docker
    \item Kubernetes
    \item Terraform
    \item Prometheus
    \item Grafana
\end{itemize}

\subsection{Success Criteria}

DevOps Engineers focus on:

\begin{itemize}
    \item Availability
    \item Scalability
    \item Security
    \item Operational efficiency
\end{itemize}

\subsection{Relationship with MLOps}

DevOps provides many foundational capabilities used by MLOps:

\begin{itemize}
    \item Automation
    \item Deployment
    \item Monitoring
    \item Infrastructure management
\end{itemize}

However, DevOps does not typically manage model lifecycle concerns such as experiment tracking or model drift.

\section{Model Validator}

\subsection{Purpose}

Model Validators provide independent assessment of machine learning models before deployment.

Their objective is to ensure that models are reliable, robust, and appropriate for their intended use.

\subsection{Responsibilities}

Typical activities include:

\begin{itemize}
    \item Independent review
    \item Methodology assessment
    \item Performance verification
    \item Documentation review
    \item Challenge and testing
\end{itemize}

\subsection{Validation Areas}

Examples include:

\begin{itemize}
    \item Data quality assessment
    \item Feature review
    \item Model assumptions
    \item Stability analysis
    \item Fairness assessment
\end{itemize}

\section{Risk Manager}

\subsection{Purpose}

Risk Managers oversee governance and regulatory compliance requirements.

They ensure that machine learning systems operate within approved risk tolerances.

\subsection{Responsibilities}

Examples include:

\begin{itemize}
    \item Governance oversight
    \item Policy enforcement
    \item Risk assessments
    \item Regulatory compliance
    \item Audit support
\end{itemize}

\subsection{Typical Frameworks}

Examples include:

\begin{itemize}
    \item SR 11-7
    \item ECB TRIM
    \item EBA Guidelines
    \item EU AI Act
\end{itemize}

\section{IT Security Officer}

\subsection{Purpose}

Security Officers ensure that machine learning platforms comply with organizational security requirements.

\subsection{Responsibilities}

Examples include:

\begin{itemize}
    \item Access control
    \item Identity management
    \item Secrets management
    \item Vulnerability management
    \item Security reviews
\end{itemize}

\subsection{Security Concerns}

Machine learning systems may expose:

\begin{itemize}
    \item Sensitive customer data
    \item Proprietary models
    \item Business-critical infrastructure
\end{itemize}

Security controls are therefore essential.

\section{Operations Team}

\subsection{Purpose}

Operations teams manage production environments after deployment.

Their objective is to maintain reliable service delivery.

\subsection{Responsibilities}

Examples include:

\begin{itemize}
    \item Incident management
    \item Service monitoring
    \item Capacity management
    \item Platform support
    \item Operational reporting
\end{itemize}

\section{Example RACI Matrix}

\subsection{Credit Scoring Project}

\begin{table}[h]
\centering
\begin{tabular}{p{5cm}cccc}
\toprule
Activity & DS & MLE & DE & RM \\
\midrule
Data Collection & C & I & R & I \\
Feature Engineering & R & C & C & I \\
Model Development & R & C & I & I \\
Validation & C & I & I & R \\
Deployment & I & R & I & A \\
Monitoring & C & R & I & A \\
Retraining & C & R & I & A \\
\bottomrule
\end{tabular}
\caption{Illustrative RACI Matrix}
\end{table}

Where:

\begin{itemize}
    \item R = Responsible
    \item A = Accountable
    \item C = Consulted
    \item I = Informed
\end{itemize}

\section{Credit Scoring End-to-End Example}

A simplified responsibility flow may be represented as:

\begin{center}
Business Owner
\end{center}

\begin{center}
$\downarrow$
\end{center}

\begin{center}
Data Engineer
\end{center}

\begin{center}
$\downarrow$
\end{center}

\begin{center}
Data Scientist
\end{center}

\begin{center}
$\downarrow$
\end{center}

\begin{center}
Model Validator
\end{center}

\begin{center}
$\downarrow$
\end{center}

\begin{center}
ML Engineer
\end{center}

\begin{center}
$\downarrow$
\end{center}

\begin{center}
Production Environment
\end{center}

\begin{center}
$\downarrow$
\end{center}

\begin{center}
Operations and Monitoring
\end{center}

Each stakeholder contributes to the successful delivery of the platform.

\section{Operating Model Best Practices}

Organizations should:

\begin{itemize}
    \item Clearly define responsibilities.
    \item Separate development and validation activities.
    \item Implement formal approval processes.
    \item Promote cross-functional collaboration.
    \item Establish governance oversight.
\end{itemize}

Well-defined responsibilities reduce operational risk and improve project success rates.

\section{Summary}

Machine learning systems require collaboration across multiple disciplines.

Key roles include:

\begin{itemize}
    \item Business Owners
    \item Product Owners
    \item Data Engineers
    \item Data Scientists
    \item ML Engineers
    \item DevOps Engineers
    \item Validators
    \item Risk Managers
    \item Security Officers
    \item Operations Teams
\end{itemize}

No single role can successfully deliver and operate enterprise machine learning systems alone.

Successful MLOps organizations establish clear responsibilities, strong collaboration mechanisms, and robust governance frameworks.

The next chapter concludes the module with a practical end-to-end MLOps project based on a Credit Scoring Platform.